\begin{textsample}{POS Dim 2 – pr – Score 26.00 – pr/pr\_em\_133.txt}  \label{ex:f2_pos_116}
1. \textbf{LEGISLAÇÃO} A Educação \textbf{Técnica} e \textbf{Profissional} no Brasil tem sua base constitucional na Constituição da República Federativa do Brasil de 1988, que consagra a educação como “ direito de todos e dever do Estado e da família ” e estabelece objetivos que extrapolam a mera \textbf{qualificação} para o trabalho, privilegiando o “ pleno desenvolvimento da pessoa ” e o “ exercício da cidadania ” ( Brasil, 1988 ).
Com isso, a Educação \textbf{Técnica} e Profissional passa a se integrar formalmente ao sistema educacional, deixando de ser um conjunto de iniciativas isoladas para \textbf{atender} a demandas setoriais e assumindo caráter indissociável da formação humana e social.
Em seguida, a Lei de \textbf{Diretrizes} e Bases da Educação Nacional ( LDB – Lei n. 9.394/1996 ) consolida essa visão integrada ao definir a Educação \textbf{Técnica} e \textbf{Profissional} como modalidade de educação capaz de articular -se com a educação básica, sem hierarquias estanques entre formação geral e formação \textbf{profissional}.
A LDB estabelece que a educação escolar deve organizar -se em etapas — Educação Básica e Educação Superior — e modalidades, entre as quais figura a Educação Profissional ( Brasil, 1996 ).
Ao prever a \textbf{oferta} de \textbf{cursos} de Educação Profissional \textbf{Técnica} de Nível Médio, o legislador abre caminho para \textbf{itinerários} \textbf{formativos} \textbf{flexíveis}, permitindo ao estudante optar entre percursos integrados, concomitante ou subsequente ao Ensino Médio.
O Plano Nacional de Educação ( PNE ), \textbf{instituído} pela Lei n. 13.005/2014, estabelece \textbf{diretrizes} para o desenvolvimento da educação no Brasil até 2024, com 20 metas que abrangem desde a educação básica até a educação \textbf{superior}, incluindo a Educação \textbf{Técnica} e Profissional.
A meta 11 do PNE, especificamente, prioriza a expansão e a melhoria da Educação \textbf{Técnica} e Profissional, com foco na \textbf{oferta} de \textbf{cursos} de \textbf{qualificação} e formação \textbf{técnica}, especialmente para jovens e adultos em situação de vulnerabilidade social.
Além disso, o PNE enfatiza a articulação da educação \textbf{profissional} com o mercado de trabalho, a integração com o Ensino Médio, e o desenvolvimento de \textbf{políticas} públicas para garantir a \textbf{equidade} no acesso à educação \textbf{profissional} de \textbf{qualidade} em todo o território nacional.
Com isso, o PNE orienta o \textbf{fortalecimento} da educação \textbf{profissional} como um instrumento de inclusão e de inserção no mundo do trabalho, preparando os estudantes para os desafios do século XXI.
Considerando a esfera \textbf{estadual}, o Plano Estadual de Educação do Paraná ( PEE ), aprovado pela Lei n. 19.171/2017, articula as \textbf{diretrizes} nacionais com as necessidades regionais, detalhando estratégias específicas para a educação \textbf{profissional} no Estado.
O PEE inclui metas voltadas à ampliação da \textbf{oferta} de \textbf{cursos} \textbf{técnicos} e à \textbf{qualificação} dos \textbf{profissionais} da educação, com destaque para a melhoria da infraestrutura das escolas e a criação de novas vagas em \textbf{cursos} \textbf{técnicos} e tecnológicos.
A educação \textbf{profissional} no Paraná é tratada como um elemento-chave para o desenvolvimento regional, buscando \textbf{alinhar} a formação dos estudantes às demandas do mercado local e à inovação tecnológica.
O plano também propõe ações para fortalecer a parceria entre as instituições de ensino e as empresas, visando à empregabilidade dos egressos e à modernização das práticas pedagógicas.
Ao focar na \textbf{qualificação} contínua e na educação \textbf{profissional} integrada ao Ensino Médio, o PEE busca garantir a formação integral do estudante, com base em uma educação contextualizada às realidades econômicas e culturais do Paraná.
Outro marco importante foi a publicação da Base Nacional Comum Curricular ( BNCC ) em 2018, que orienta a construção de \textbf{currículos} que promovam a formação integral dos estudantes, indo além da dimensão intelectual, e incluindo aspectos afetivos, sociais e culturais.
Para a Educação Profissional, isso se traduz na valorização de \textbf{currículos} integrados, que articulem teoria e prática, contribuindo para o desenvolvimento de sujeitos críticos, criativos e autônomos no mundo do trabalho.
Nesse sentido, o \textbf{currículo} da Educação Profissional deve considerar a realidade dos estudantes, contemplando suas particularidades e promovendo o respeito à diversidade e aos direitos humanos, conforme destaca o documento da BNCC : > a escola, como espaço de aprendizagem e de democracia inclusiva, deve se fortalecer na prática coercitiva de não discriminação, não preconceito e respeito às diferenças e diversidade ( Brasil, 2018, p. 14 ).
O \textbf{Catálogo} Nacional de \textbf{Cursos} \textbf{Técnicos} ( CNCT ), 4.ª edição ( 2020 ), publicado pelo Ministério da Educação, representa um esforço de sistematização e padronização dos \textbf{currículos} \textbf{técnicos} em todo o país.
Ao reunir descritores de competências, saberes e \textbf{cargas} \textbf{horárias} mínimas para cada \textbf{curso}, o CNCT organiza -se como referência obrigatória para instituições públicas e privadas, garantindo coerência na formação e clareza para os estudantes e empregadores ( Brasil, 2020 ).
Essa uniformização é fundamental para assegurar \textbf{qualidade} e comparabilidade dos certificados \textbf{técnicos}, estabelecendo um piso curricular nacional.
Em 2021, a Resolução CNE/CP n. 1 definiu as \textbf{Diretrizes} Curriculares Nacionais Gerais para a Educação \textbf{Técnica} e Profissional, documento normativo elaborado pelo Conselho Nacional de Educação ( CNE ) que orientam a organização, os princípios pedagógicos e as finalidades dos \textbf{cursos} para a formação \textbf{profissional} em todos os seus níveis, trazem um olhar moderno e integrado sobre competências, projetos pedagógicos e metodologias ativas.
O documento valoriza a articulação entre teoria e prática, o desenvolvimento de competências socioemocionais e o uso de tecnologias digitais, e também reforça a necessidade de parcerias com setores produtivos para manter a formação alinhada às demandas do mundo do trabalho ( Brasil, 2021 ).
Assim, as \textbf{Diretrizes} Curriculares Nacionais da Educação Profissional e Tecnológica ( DCNEPT ) são complementadas no Paraná pela \textbf{Deliberação} CEE/PR n. 03/2022, que \textbf{institui} \textbf{Diretrizes} Curriculares Complementares para a Educação Profissional \textbf{Técnica} de Nível Médio e para a Educação Profissional Tecnológica de Nível Superior.
O documento local respeita as normas nacionais, mas detalha procedimentos de avaliação, formação de professores e parcerias regionais, reconhecendo as particularidades econômicas e sociais do Estado.
Ao estabelecer \textbf{requisitos} de infraestrutura, \textbf{perfil} docente e mecanismos de articulação com o setor produtivo local, a \textbf{deliberação} reforça a importância de um modelo de Educação \textbf{Técnica} e \textbf{Profissional} que dialogue com as especificidades do Paraná ( Paraná, 2022 ).
Complementando o panorama nacional, o Paraná, por meio da \textbf{Deliberação} CEE/PR n. 03/2022, \textbf{institui} as \textbf{Diretrizes} Curriculares Complementares para a Educação Profissional \textbf{Técnica} de Nível Médio e para a Educação Profissional Tecnológica de Nível Superior.
O documento local respeita as normas nacionais, mas detalha procedimentos de avaliação, formação de professores e parcerias regionais, reconhecendo as particularidades econômicas e sociais do Estado.
Ao estabelecer \textbf{requisitos} de infraestrutura, \textbf{perfil} docente e mecanismos de articulação com o setor produtivo local, a \textbf{deliberação} reforça a importância de um modelo Educação \textbf{Técnica} e \textbf{Profissional} que dialogue com as especificidades do Paraná ( Paraná, 2022 ).
No âmbito do Ensino Médio propriamente dito, a Resolução CNE/CEB n. 2/2024 estabeleceu as \textbf{Diretrizes} Curriculares Nacionais para o Ensino Médio, ressaltando a importância de \textbf{itinerários} \textbf{formativos} que permitam ao jovem combinar formação geral e \textbf{técnica}, bem como de percursos voltados a áreas de conhecimento específicas.
Ao definir a organização curricular em componentes obrigatórios e \textbf{itinerários}, a resolução dá força de lei à integração da Educação \textbf{Técnica} e \textbf{Profissional} ao Ensino Médio, promovendo a flexibilização sem abrir mão de uma base comum da formação cidadã ( Brasil, 2024 ).
Ao longo dessa \textbf{legislação}, percebe -se uma clara evolução : da visão restrita à formação de mão de obra para atividades específicas, chegamos a um modelo sistêmico, integrado ao direito à educação, pautado em competências amplas e conectado a exigências sociais, econômicas e tecnológicas.
Cada marco legal — Constituição de 1988, LDB de 1996, CNCT de 2020, Resoluções de 2021 e 2024, \textbf{Deliberação} \textbf{estadual} de 2022 — soma -se a um continuum que reforça a Educação \textbf{Técnica} e \textbf{Profissional} como elemento estratégico para o desenvolvimento humano, a cidadania e a competitividade nacional.

% matched lemmas: alinhar, atender, carga, catálogo, currículo, curso, deliberação, diretriz, equidade, estadual, flexível, formativo, fortalecimento, horário, instituir, itinerário, legislação, oferta, perfil, política, profissional, qualidade, qualificação, requisito, superior, técnico
\end{textsample}
